% ==============================================================================
% ch:uvod — úvod práce: kontext, motivace, struktura
% Bez citací (odkazy pouze na vlastní kapitoly).
% ==============================================================================

Český trh práce vstupuje do~období, v~němž se překrývají technologická změna, demografické stárnutí a~vyčerpávání modelu levnější práce.
Automatizace, robotizace a~generativní \acs{AI} mění poptávku po~kvalifikaci a~rychlost přechodů mezi profesemi; růst cenové a~mzdové hladiny ve~střední Evropě současně zmenšuje prostor pro konkurenceschopnost založenou převážně na~nízké jednotkové ceně práce.
Česká republika si udržuje vysokou zaměstnanost a~relativně nízkou příjmovou nerovnost, ale tyto příznivé makroekonomické indikátory mohou zakrývat slabší kupní sílu hodinového příjmu, delší pracovní dobu a~nízkou schopnost kolektivně vyjednat podmínky nastávající strukturální změny.\par

Hlavní analytická linie práce vychází z~tohoto napětí. Nízká odborová organizovanost, nízké pokrytí \aclins{KS} a~slabé zapojení sociálních partnerů do~\aclgen{APZ} nejsou jen tři oddělené statistiky; společně ukazují, zda se náklady technologických a~odvětvových přesunů rozloží mezi sociální partnery, nebo zda zůstanou převážně na~jednotlivých pracovnících. Sociální dialog je proto v~této práci chápán jako nástroj řízení strukturální změny: vyjednává mzdové standardy, převádí produktivitu do~příjmů, omezuje informační asymetrii a~může spoluutvářet rekvalifikační infrastrukturu.\par

Toto téma je zároveň bilanční otázkou české evropské konvergence. Po~více než dvou desetiletích členství v~\aclloc{EU} již nestačí hodnotit úspěch pouze podle dohánění \aclgen{HDP} nebo podle zapojení do~dodavatelských řetězců. Podstatné je, zda se vytvořená hodnota promítá také do~pracovních standardů, vyjednávací síly zaměstnanců a~schopnosti přejít k~vyšší přidané hodnotě. Slabý sociální dialog v~tomto rámci není okrajovým tématem pracovního práva, nýbrž jedním z~indikátorů, zda se \acl{ČR} stabilizuje jako levnější subdodavatelská periferie, nebo zda dokáže udržitelnými sociálními podmínkami a~kvalifikovanou prací podpořit kvalitativní růst.\par

Práce proto ověřuje hypotézu: \emph{Posílení sociálního dialogu a~kolektivního vyjednávání v~\acloc{ČR} by přispělo k~udržitelnému hospodářskému rozvoji.} Cílem není prokázat jednoduchou kauzalitu mezi jedním právním institutem a~jedním ekonomickým výsledkem, ale posoudit napříč deskriptivními daty, evropskou komparací, korelační analýzou a~případovými studiemi, zda slabý sociální dialog souvisí s~horšími výsledky trhu práce a~zda by jeho posílení podpořilo udržitelný hospodářský rozvoj.\par

Volba referenčních zemí pro komparativní část staví \acacc{ČR} do~několika typologicky odlišných srovnání:
\begin{itemize}
  \item \acl{DK} v~současnosti představuje ekonomiku s~nejrychleji rostoucím \acs{HDP} v~\acs{EU} s~modelem flexicurity, v~němž je flexibilita spojena se silnou \aclins{APZ} a~tradicí kolektivního vyjednávání a~vysokou odborovou organizovaností,
  \item \acl{AT} ukazuje, jak lze i~při nižší odborové organizovanosti dosáhnout vysokého pokrytí \aclpgen{KS} prostřednictvím odvětvové koordinace a~rozšiřování závaznosti \aclpgen{KSVS},
  \item \acl{DE} jako hlavní obchodní partner \aclgen{ČR} nabízí varovnou zkušenost s~deregulací nízkopříjmové práce,
  \item Polsko a~Slovensko poskytují postkomunistické srovnání, v~němž podobná historická výchozí pozice vedla k~podobným trajektoriím vývoje trhu práce a~odborového hnutí.
\end{itemize}

Struktura práce sleduje stejnou argumentační logiku. Teoretická část vymezuje historický a~legislativní rámec, pojmy \acl{KV}, \acl{KS}, sociální partneři, \acl{APZ} a~nové formy práce a~zasazuje je do~současné literatury. Metodologická kapitola popisuje datové zdroje, přepočet kupní síly a~modelové výpočty. Empirická část nejprve analyzuje český trh práce, pokrytí \aclpins{KS}, \aclacc{APZ}, výhledové faktory a~evropské modely, poté výsledky koncentruje v~korelační analýze, případových studiích a~SWOT analýze. Návrhová kapitola nakonec převádí zjištěné slabiny do~pěti inovací: aktivace rozšiřování \aclpgen{KSVS}, zapojení sociálních partnerů do~\aclgen{APZ}, mzdové transparentnosti, vyrovnání postavení ekonomicky závislé práce a~veřejného digitálního registru \aclpgen{KS}.\par
